% Existing source excerpts, 2026-10-11. Attribution and modification log: ../README.md.
% Onion12138/math excerpts and associated original figures: LPPL 1.3c.
% higyg/middleschoolmath-ebook excerpt and associated figures: GPL v3.
% Extracted without editing their problem conditions or values; not a standalone LaTeX document.

% Upstream Onion12138/math, test.tex; source commit and links in README.
如图，正方形$ABCD$以$AB$为边向外作等边三角形$ABE$，$M$为对角线$BD$上一点，将$BM$绕$B$逆时针旋转$60^\circ$到$BN$，若正方形边长为2，求
    $AM+BM+CM$的最小值。
\end{example}

% Upstream Onion12138/math, test.tex; source commit and links in README.
如图，$C,D$为线段$AB$上两点，$P$是线段$CD$上的动点，分别以$AP,BP$为边在$AB$同侧作两个等边三角形$APE$和$BPF$，$M$是$EF$的中点，
    已知$AB=20,AC=BD=2$，当$P$从$C$运动到$D$时(无重复运动)，求$M$点的运动路径长。
\end{problem}

% Upstream Onion12138/math, test.tex; source commit and links in README.
(1)如图1，在正方形$ABCD$中，$E,F$为$BC,CD$上点，且$\angle EAF=45^\circ$，直接写出$BE,DF,EF$三条线段的数量关系。 \\
    (2)如图2，将正方形$ABCD$改为四边形$ABCD$，$AB=AD,\angle B+\angle D=180^\circ$，$E,F$分别为边$BC,CD$上的点，
    $\angle EAF=\frac{1}{2}\angle BAD$，判断(1)中结论是否依然成立。 \\
    (3)在(2)基础上将$\triangle AEF$绕$A$逆时针旋转，$E,F$运动到$BC,CD$延长线上，如图3，判断$BE,DF,EF$三条线段的数量关系并证明。 \\
\end{problem}

% Upstream higyg/middleschoolmath-ebook, ItemBankBody/T22-01.tex, T22-A01-10.
\begin{defproblem}{T22-A01-10}%
\begin{onlyproblem}%
% \href{run:./ItemBankFigures/T22-A01-06.ggb}{（查看动态演示）}
（2017年河南中考）
如图1，在Rt$\triangle ABC$中，$\angle  A =90 ^{\circ }$，$AB=AC$，点$D$，$E$分别在边$AB$，$AC$上，$AD=AE$，连接$DC$，点$M$，$P$，$N$分别为$DE$，$DC$，$BC$的中点．

（1）观察猜想

图1中，线段$PM$与$PN$的数量关系是{\_}{\_}{\_}{\_}{\_}{\_}{\_}{\_}{\_}{\_}{\_}，位置关系是{\_}{\_}{\_}{\_}{\_}{\_}{\_}{\_}{\_}；

（2）探究证明

把$\triangle ADE$绕点$A$逆时针方向旋转到图2的位置，连接$MN$，$BD$，$CE$，判断$\triangle PMN$的形状，并说明理由；

（3）拓展延伸

把$\triangle ADE$绕点$A$在平面内自由旋转，若$AD=4$，$AB=10$，请直接写出$\triangle PMN$面积的最大值．


\begin{center}
\includegraphics[scale=0.8]{T22-A01-10-01.pdf}\qquad
\includegraphics[scale=0.8]{T22-A01-10-02.pdf}
\end{center}


% \begin{flushright}
% \includegraphics[scale=0.8]{T22-A01-02-01.pdf}\\
% \includegraphics[scale=0.8]{T22-A01-02-02.pdf}\\
% \includegraphics[scale=0.8]{T22-A01-02-03.pdf}
% \end{flushright}

\end{onlyproblem}%
\begin{onlysolution}%
（1）$PM=PN$；$PM \bot  PN$；

（2）$\triangle PMN$为等腰直角三角形，理由略；

（3）$\triangle PMN$面积的最大值为$\dfrac{49}{2}$．
\begin{center}
% \includegraphics[width=.95\textwidth]{T22-A01-06-04-Ans.jpg}
\end{center}
\end{onlysolution}%
\end{defproblem}
